\section{Native NLP and Interior-Point Details}
\label{sec:engine-nlp-ipm}

The nonlinear OPF solvers are local to \texttt{src/optimal\_power\_flow}.  They
do not dispatch through an external NLP engine by default.  The AC OPF entry
point selects either a native reduced-space primal-dual path or the parity
full-space IPM path described in Section~\ref{sec:opf}.

\subsection{Native AC OPF Loop}

The native path in \texttt{ac\_opf.cpp} uses variables for non-slack AC angles,
PQ voltage magnitudes, generator dispatch, DC voltage magnitudes, and VSC
converter powers.  Equality rows cover AC balances, DC balances, and VSC energy
balance.  Bound constraints are handled with barrier terms, while nonlinear
branch and converter limits enter through exterior penalty terms.  If this path
fails and \texttt{ACOPFOptions::allow\_fallback} is true, the code falls back to
the fast dispatch plus AC PF projection route.

\subsection{Parity IPM Loop}

The parity path builds explicit equality and inequality rows in
\texttt{parity\_formulation.cpp} and solves the resulting full-space NLP in
\texttt{parity\_ipm.cpp}.  The solve loop uses a Mehrotra predictor-corrector
step with second-order complementarity correction, iterative refinement, and
fraction-to-boundary primal/dual step lengths.  Current centering is
\begin{align}
\sigma = \left(\frac{\mu_{aff}}{\mu_{cur}+10^{-16}}\right)^3
\end{align}
after clamping the ratio to \([0,1]\).  The implementation also applies up to
two extra Gondzio corrector solves when enabled by the current loop logic.

The KKT system is solved with Eigen \texttt{SparseLU} on sparse systems or
\texttt{PartialPivLU} on dense systems.  There is no separate local NLP/IPM
engine source tree beneath \texttt{src/engine}.
